\section{Conclusion}
We have constructed continuous curves whose real-axis contacts recover prime
values and twin-prime starts. The proofs connect exact integer samples of
normalized Fejér waves with the divisor function. Activation makes the
infinite construction locally finite, while nonnegative heights make a joint
zero an exact simultaneous survival condition.

For integer $m\ge2$ and integer $n$, the kernel bounds and samples are
\begin{equation*}
\begin{aligned}
0&\le W_m(t)\le1
&&\text{[Wave Bounds]}\\
W_m(n)&=\mathbf 1_{m\mid n}.
&&\text{[Divisibility Samples]}
\end{aligned}
\end{equation*}

For $M\ge2$ and $t\ge2$, the activated update preserves the real coordinate
and raises only imaginary height:
\begin{equation*}
\begin{aligned}
F_M(t)&=F_{M-1}(t)+i a_M(t)W_M(t)
&&\text{[Activated Recurrence]}\\
h_M(t)&\ge h_{M-1}(t).
&&\text{[Monotone Height]}
\end{aligned}
\end{equation*}

The added term is zero on $[2,M]$. On any interval $[2,X]$, $X\ge2$,
sufficiently late stages agree with the continuous limit exactly:
\begin{equation*}
\begin{aligned}
M\ge\lceil X\rceil&\ \Longrightarrow\ F_M(t)=F(t)\quad(2\le t\le X).
&&\text{[Compact Stabilization]}
\end{aligned}
\end{equation*}

At integers $n\ge2$ the height is the classical divisor count minus two;
noninteger exclusion then identifies every real-axis contact on $t\ge2$:
\begin{equation*}
\begin{aligned}
h(n)&=\tau(n)-2
&&\text{[Divisor Height]}\\
F(t)\in\mathbb R&\ \Longleftrightarrow\ t\text{ is prime}.
&&\text{[Prime Contacts]}
\end{aligned}
\end{equation*}

For composite $M$ and integer $n\ge2$, a filter can add divisor information
without changing contact membership:
\begin{equation*}
\begin{aligned}
h_M(n)=0&\ \Longleftrightarrow\ h_{M-1}(n)=0.
&&\text{[Composite Redundancy]}
\end{aligned}
\end{equation*}

At prime step $p$, the first newly rejected integer is $p^2$:
\begin{equation*}
\begin{aligned}
h_{p-1}(p^2)&=0,\quad h_p(p^2)=1.
&&\text{[Prime-Square Rejection]}
\end{aligned}
\end{equation*}

For every integer $M\ge1$ and integer $n\ge2$, the finite-stage certificate is
\begin{equation*}
\begin{aligned}
2\le n < (M+1)^2,\ h_M(n)=0
&\ \Longrightarrow\ n\text{ is prime}.
&&\text{[Finite Certification]}
\end{aligned}
\end{equation*}

If $q$ is the least prime greater than $M\ge1$, the sharper boundary is
\begin{equation*}
\begin{aligned}
n\text{ composite},\ 2\le n < q^2
&\ \Longrightarrow\ h_M(n) > 0
&&\text{[Sharp Certification]}\\
h_M(q^2)&=0.
&&\text{[First Composite Contact]}
\end{aligned}
\end{equation*}

For $n=\prod_jp_j^{e_j}\ge2$, the divisor-product formula yields the low-height
classifications and square-parity law:
\begin{equation*}
\begin{aligned}
h(n)+2&=\prod_j(e_j+1)
&&\text{[Divisor Factorization]}\\
h(n)\text{ is odd}&\ \Longleftrightarrow\ n\text{ is a square}.
&&\text{[Square Parity]}
\end{aligned}
\end{equation*}

The low-height forms, for integer $n\ge2$ and suitable primes $p,q$, are
\begin{equation*}
\begin{aligned}
h(n)=0&\ \Longleftrightarrow\ n=p
&&\text{[Prime Height]}\\
h(n)=1&\ \Longleftrightarrow\ n=p^2
&&\text{[Prime-Square Height]}\\
h(n)=2&\ \Longleftrightarrow\ n=p^3\text{ or }n=pq\ (p\ne q)
&&\text{[Height-Two Types]}\\
h(n)=3&\ \Longleftrightarrow\ n=p^4.
&&\text{[Prime-Fourth-Power Height]}
\end{aligned}
\end{equation*}

For coprime integers $a,b\ge2$, independent divisor choices give
\begin{equation*}
\begin{aligned}
h(ab)+2&=(h(a)+2)(h(b)+2).
&&\text{[Coprime Multiplication]}
\end{aligned}
\end{equation*}

Applying the same waves at both endpoints gives the pair height. For $t\ge2$
and integer $n\ge2$,
\begin{equation*}
\begin{aligned}
B_M(t)&=2\sin^2(\pi t)+\sum_{m=2}^{M}V_m(t)
&&\text{[Pair-Wave Construction]}\\
B(n)&=\tau(n)+\tau(n+2)-4.
&&\text{[Pair Divisor Height]}
\end{aligned}
\end{equation*}
\begin{equation*}
\begin{aligned}
T(t)\in\mathbb R
&\ \Longleftrightarrow\ t\text{ and }t+2\text{ are prime}.
&&\text{[Twin Contacts]}
\end{aligned}
\end{equation*}

The same finite-stage conclusion holds in the certified range
$n+2 < (M+1)^2$. Bare-wave folding holds for integer $m\ge2$, $k\ge1$ and
all real $t$:
\begin{equation*}
\begin{aligned}
\sum_{j=0}^{k-1}W_{km}(t+jm)&=W_m(t).
&&\text{[Translated Folding]}
\end{aligned}
\end{equation*}

For the activated filter-$4$ pair on $t\ge2$, the correction preserves exact
initial values and disappears throughout $t\ge5$:
\begin{equation*}
\begin{aligned}
V_4(t)&=W_2(t)-C_4(t)
&&\text{[Activation Correction]}\\
t\ge5&\ \Longrightarrow\ C_4(t)=0.
&&\text{[Tail Identity]}
\end{aligned}
\end{equation*}

Finally, the complementary product gives the same twin contacts with bounded
imaginary height:
\begin{equation*}
\begin{aligned}
0&\le A(t)\le1
&&\text{[Bounded Amplitude]}\\
\widetilde T(t)\in\mathbb R&\ \Longleftrightarrow\ A(t)=1
&&\text{[Unit-Level Contacts]}\\
A(t)=1&\ \Longleftrightarrow\ t,t+2\text{ are prime},\quad t\ge2.
&&\text{[Twin Equivalence]}
\end{aligned}
\end{equation*}

Together, these properties characterize the contact sets and their finite
filter updates in two continuous representations. They also identify the
precise arithmetic meaning of the additive height. The location and
frequency of simultaneous contacts remain the quantitative question, and no
twin-prime infinitude theorem follows from these representations alone.
